\setcounter{section}{-1}

\section{Preliminaries}

\subsection{Basics}

For Exercises 1 to 4, let
\[
M = 
\begin{pmatrix}
    1 & 1 \\
    0 & 1
\end{pmatrix}
\]
and $\bb = \{X \in \ac \mid MX = XM\}$, where $\ac$ denotes the set of $2 \times 2$ matrices with real entries.

\begin{exercise}\label{ex0.1.1}
    Determine which of the following elements of $\ac$ lie in $\bb$:
    \[
    \begin{pmatrix}
        1 & 1 \\
        0 & 1
    \end{pmatrix},
    \begin{pmatrix}
        1 & 1 \\
        1 & 1
    \end{pmatrix},
    \begin{pmatrix}
        0 & 0 \\
        0 & 0
    \end{pmatrix},
    \begin{pmatrix}
        1 & 1 \\
        1 & 0
    \end{pmatrix},
    \begin{pmatrix}
        1 & 0 \\
        0 & 1
    \end{pmatrix},
    \begin{pmatrix}
        0 & 1 \\
        1 & 0 
    \end{pmatrix}
    \]
\end{exercise}

\begin{sol}
    The 1st, 3rd, 5th, and 6th matrices are $M$, the zero matrix, the identity matrix, and the exchange matrix respectively. Among these, the 1st, 3rd, and 5th matrices belong to $\bb$. Checking the 2nd matrix:
    \begin{align*}
        MX & = 
        \begin{pmatrix}
            1 & 1 \\
            0 & 1
        \end{pmatrix}
        \begin{pmatrix}
            1 & 1 \\
            1 & 1
        \end{pmatrix} = 
        \begin{pmatrix}
            1 \cdot 1 + 1 \cdot 1 & 1 \cdot 1 + 1 \cdot 1 \\
            0 \cdot 1 + 1 \cdot 1 & 0 \cdot 1 + 1 \cdot 1
        \end{pmatrix} = 
        \begin{pmatrix}
            2 & 2 \\ 1 & 1
        \end{pmatrix} \\
        XM & = 
        \begin{pmatrix}
            1 & 1 \\
            1 & 1 
        \end{pmatrix}
        \begin{pmatrix}
            1 & 1 \\
            0 & 1
        \end{pmatrix} = 
        \begin{pmatrix}
            1 \cdot 1 + 1 \cdot 0 & 1 \cdot 1 + 1 \cdot 1 \\
            1 \cdot 1 + 1 \cdot 0 & 1 \cdot 1 + 1 \cdot 1 
        \end{pmatrix} =
        \begin{pmatrix}
            2 & 1 \\
            2 & 1
        \end{pmatrix}
    \end{align*}
    Then $MX \neq XM$ and the 2nd matrix does not belong to $\bb$. Similarly, the 4th matrix does not belong to $\bb$.
\end{sol}

\begin{exercise}\label{ex0.1.2}
    Prove that if $P, Q \in \bb$, then $P + Q \in \bb$ where $+$ denotes the usual sum of two matrices.
\end{exercise}

\begin{sol}
    We calculate the following:
    \begin{align*}
        M(P + Q) & = 
        \begin{pmatrix}
            1 & 1 \\
            0 & 1
        \end{pmatrix}
        \begin{pmatrix}
            a + e & b + f \\
            c + g & d + h
        \end{pmatrix} \\
        & = 
        \begin{pmatrix}
            a + e + c + g & b + f + d + h \\
            c + g & d + h
        \end{pmatrix} \\
        & = 
        \begin{pmatrix}
            a + c & b + d \\
            c & d
        \end{pmatrix} + 
        \begin{pmatrix}
            e + g & f + h \\
            g & h
        \end{pmatrix} \\
        & = MP + MQ \\
        & = PM + QM \\
        & = 
        \begin{pmatrix}
            a & a + b \\
            c & c + d
        \end{pmatrix} + 
        \begin{pmatrix}
            e & e + f \\
            g & g + h
        \end{pmatrix} \\
        & = 
        \begin{pmatrix}
            a + e & a + e + b + f \\
            c + g & c + g + d + h
        \end{pmatrix} \\
        & = 
        \begin{pmatrix}
            a + e & b + f \\
            c + g & d + h
        \end{pmatrix}
        \begin{pmatrix}
            1 & 1 \\
            0 & 1
        \end{pmatrix} \\
        & = (P + Q)M \qh
    \end{align*}
\end{sol}

\begin{exercise}\label{ex0.1.3}
    Prove that if $P, Q \in \bb$, then $P\cdot Q \in \bb$ where $\cdot$ denotes the usual product of two matrices.
\end{exercise}

\begin{sol}
    Use a similar process as in \hyperref[ex0.1.2]{Exercise 0.1.2}.
\end{sol}

\begin{exercise}\label{ex0.1.4}
    Find conditions on $p, q, r, s$ which determine precisely when
    \[
    \begin{pmatrix}
        p & q \\
        r & s
    \end{pmatrix} \in \bb.
    \]
\end{exercise}

\begin{sol}
    Let $X$ be the matrix described above. Note that
    \[
    MX =
    \begin{pmatrix}
        p + r & q + s \\
        r & s
    \end{pmatrix} \quad
    XM = 
    \begin{pmatrix}
        p & p + q \\
        r & r + s
    \end{pmatrix}
    \]
    Because $X \in \bb$, then we may compare entries to obtain the following:
    \[
    \begin{cases}
        p + r = p \\
        q + s = p + q \\
        r = r \\
        s = r + s
    \end{cases}
    \]
    The first and fourth equations force $r = 0$, and the second equation forces $p = s$. Then $\bb$ is classified as
    \[\bb = \set*{\left.
    \begin{pmatrix}
        p & p + q \\
        0 & p
    \end{pmatrix}\,\right|\, p, q \in \r} \qh\]
\end{sol}

\begin{exercise}\label{ex0.1.5}
    Determine whether the following functions $f$ are well-defined:
    \begin{subproblems}
        \item $f : \q \to \z$ defined by $f(a/b) = a$.
        \item $f : \q \to \q$ defined by $f(a/b) = a^2/b^2$.
    \end{subproblems}
\end{exercise}

\begin{solalph}
    \item No, because $1/2 = 2/4$, but $f(1/2) = 1$ and $f(2/4) = 2$.
    \item Yes; suppose $a/b = c/d$. Then $ad = bc$, or that $a^2d^2 = b^2c^2$. Then $a^2/b^2 = c^2/d^2$, or $f(a/b) = f(c/d)$. 
\end{solalph}

\begin{exercise}\label{ex0.1.6}
    Determine whether the function $f: \r^+ \to \z$ defined by mapping a real number $r$ to the first digit to the right of the decimal point in a decimal expansion of $r$ is well-defined.
\end{exercise}

\begin{sol}
    No. Note that $1=1.000\ldots = 0.999\ldots$, but $f(1.000\ldots) = 1$ and $f(0.999\ldots) = 9$.
\end{sol}

\begin{exercise}\label{ex0.1.7}
    Let $f: A \to B$ be a surjective map of sets. Prove that the relation
    \[a \sim b \iff f(a) = f(b)\]
    is an equivalence relation whose equivalence classes are the fibers of $f$.
\end{exercise}

\begin{sol}
    Since $=$ is an equivalence relation on $B$, then $\sim$ is an equivalence relation.

    Consider an equivalence class of some $a \in A$, which is the set $\set{x \in A \mid x \sim a}$. By definition of $\sim$, this is the set $\set{x \in A \mid f(x) = f(a)}$, which is precisely the fiber of $f$ over $f(a)$. Since $f$ is surjective, every fiber of $f$ is nonempty, and every fiber corresponds to some equivalence class of $\sim$.
\end{sol}

\newpage

\subsection{Properties of the Integers}

\begin{exercise} \label{ex0.2.1}
    For each of the following pairs of integers $a$ and $b$, determine their greatest common divisor, their least common multiple, and write their greatest common divisor in the form $ax + by$ for some integers $x$ and $y$.
    \begin{subproblems}
        \item $a = 20, b = 13$
        \item $a= 69, b = 372$
        \item $a = 792, b = 275$
        \item $a = 11391, b = 5673$
        \item $a= 1761, b = 1567$
        \item $a = 507885, b = 60808$
    \end{subproblems}
\end{exercise}

\begin{sol}
    For this exercise, we will only do (e) as that has the most steps in calculating both the $\gcd$ and the Euclidean Algorithm. The $\lcm$ is obtained by dividing the product $ab$ by $(a, b)$.
    \begin{subproblems}
        \item$(20, 13) = 1, \lcm(20, 13) = 260, 1 = 2(20) - 3(13)$
        \item $(69, 372) = 3,\lcm(69, 372) = 8556, 27(69) -5(372)$
        \item $(792, 275) = 11, \lcm(792, 275) = 19800, 8(792) - 23(275)$
        \item $(11391, 5673) = 3 ,\lcm(11391, 5673) = 21540381, 3 = 253(5673) - 126(11391)$
        \item Applying the Euclidean Algorithm to $a = 1761$ and $b = 1567$, we get:
        \begin{align*}
            1761 & = (1)1567 + 194 \\
            1567 & = (8)194 + 15 \\
            194 & = (12)15 + 14 \\
            15 & = (1)14 + 1
        \end{align*}
        Then $(1761, 1567) = 1$ so that $\lcm(1761, 1567) = 2759487$. Reversing the Euclidean Algorithm steps to solve for 1, we get:
        \begin{align*}
            1 & = 15 - 14 \\
            & = 15 - (194 - 12(15)) = 13(15) - 194 \\
            & = 13(1567 - 8(194)) - 194 = 13(1567) - 105(194) \\
            & = 13(1567) - 105(1761 - 1567) = -105(1761) + 118(1567)
        \end{align*}
        \item $(507885, 60808) = 691, \lcm(507885, 60808) = 44693880, 691 = 142(60808) - 17(507885)$ \qh
    \end{subproblems}
\end{sol}

\begin{exercise}\label{ex0.2.2}
    Prove that if the integer $k$ divides the integers $a$ and $b$, then $k$ divides $as + bt$ for every pair of integers $s$ and $t$. 
\end{exercise}

\begin{sol}
    Since $k \mid a$ and $k \mid b$, there exists $x, y \in \z$ such that $a = kx$ and $b = ky$. Then for any $s, t \in \z$, we have $as + bt = kxs + kyt = k(xs + yt)$ which is divisible by $k$.
\end{sol}

\begin{exercise}\label{ex0.2.3}
    Prove that if $n$ is composite, then there are integers $a$ and $b$ such that $n$ divides $ab$, but $n$ does not divide either $a$ or $b$.
\end{exercise}

\begin{sol}
    By the Fundamental Theorem of Arithmetic, $n$ has at least two prime factors $a, b$ such that $1 < a, b < n$. Putting $n = ab$, then $n \mid ab$ but $n \nmid a$ and $n \nmid b$.
\end{sol}

\begin{exercise}\label{ex0.2.4}
    Let $a, b$ and $N$ be fixed integers with $a$ and $b$ nonzero, and let $d = (a, b)$ be the greatest common divisor of $a$ and $b$. Suppose $x_0$ and $y_0$ are particular solutions to $ax + by = N$. Prove, for any integer $t$, that the integers
    \[x = x_0 + \frac bdt \quad \text{and} \quad y = y_0 - \frac adt\]
    are also solutions to $ax + by = N$.

    \remark{This is, in fact, the general solution.}
\end{exercise}

\begin{sol}
    \begin{align*}
        ax + by & = a\left(x_0 + \frac bdt\right) + b\left(y_0 - \frac adt\right) \\
        & = ax_0 + \frac{ab}{d}t + by_0 - \frac{ab}{d}t \\
        & = ax_0 + by_0 = N \qh
    \end{align*}
\end{sol}

\begin{exercise}\label{ex0.2.5}
    Determine the value $\phi(n)$ for each integer $n \leq 30$, where $\phi$ denotes the Euler $\phi$-function.
\end{exercise}

\begin{sol}
    \[\begin{array}{c|c|c|c|c|c|c|c|c|c|c|c|c|c|c|c}
        n & 1 & 2 & 3 & 4 & 5 & 6 & 7 & 8 & 9 & 10 & 11 & 12 & 13 & 14 & 15 \\
        \hline
        \phi(n) & 1 & 1 & 2 & 2 & 4 & 2 & 6 & 4 & 6 & 4 & 10 & 4 & 12 & 6 & 8 \\
    \end{array}\]
    \[\begin{array}{c|c|c|c|c|c|c|c|c|c|c|c|c|c|c|c}
        n & 16 & 17 & 18 & 19 & 20 & 21 & 22 & 23 & 24 & 25 & 26 & 27 & 28 & 29 & 30 \\
        \hline
        \phi(n) & 8 & 16 & 6 & 18 & 8 & 12 & 10 & 22 & 8 & 20 & 12 & 18 & 12 & 28 & 8 \\
    \end{array}\]
\end{sol}

\begin{exercise}\label{ex0.2.6}
    Prove the Well Ordering Property of $\z$ by induction, and prove the minimal element is unique.
\end{exercise}

\begin{sol}
    Let $S \subseteq \zp$ be nonempty. We proceed by induction that $S$ has a minimal element. Assume, by way of contradiction, that $S$ has no minimal element, and let $S'$ be the set of elements that are not in $S$. Since the minimal element of $\zp$ is 0, then $0 \not\in S$ so that $0 \in S'$. Now for some $k \in \zp$, suppose every integer $j$ such that $0 \leq j \leq k$ is in $S'$. Then $k + 1 \not\in S$, since if it were, then it would be the minimal element of $S$ (as every integer less than or equal to $k$ is in $S'$). Thus, $k + 1 \in S'$. By induction, every integer in $\zp$ is in $S'$, contradicting that $S$ is nonempty. Thus, $S$ has a minimal element.

    To prove uniqueness of the minimal element, suppose $S$ has two minimal elements $x$ and $y$. Then by definition of minimality, $x \leq y$ and $y \leq x$, so that $x = y$.
\end{sol}

\begin{exercise}\label{ex0.2.7}
    If $p$ is a prime, prove that there do not exist nonzero integers $a$ and $b$ such that $a^2 = pb^2$.
\end{exercise}

\begin{sol}
    Assume, by way of contradiction, that there exist nonzero integers $a$ and $b$ such that $a^2 = pb^2$. Without loss of generality, we may assume that $a$ and $b$ have no common factors (otherwise, we could divide them both by their greatest common divisor). Then $p \mid a^2$, so that $p \mid a$. Then there exists some integer $k$ such that $a = pk$. We then have
    \[a^2 = (pk)^2 = p^2k^2 = pb^2\]
    so that $b^2 = pk^2$. Then $p \mid b^2$, so that $p \mid b$. However, this contradicts our assumption that $a$ and $b$ have no common factors. Thus, there do not exist nonzero integers $a$ and $b$ such that $a^2 = pb^2$.
\end{sol}

\begin{exercise} \label{ex0.2.8}
    Let $p$ be a prime, $n \in \zp$. Find a formula for the largest power of $p$ which divides $n! = n(n - 1)(n - 2) \ldots 2 \cdot 1$.

    \remark{It involves the greatest integer function.}
\end{exercise}

\begin{sol}
    Note that there exists $k_1 \in \zp$ such that $n \geq k_1p$. Then the integers $p, 2p, \ldots, k_1p$ contribute $k_1$ factors of $p$ in $n!$. Next, there exists $k_2 \in \zp$ such that $n \geq k_2p^2$. Then the integers $p^2, 2p^2, \ldots, k_2p^2$ contribute $k_2$ additional factors of $p$ in $n!$. We continue this process until some $k_m \in \zp$ such that $n \geq k_m p^m$ but $n < k_{m+1} p^{m+1}$. The largest power of $p$ which divides $n!$ is then given by
    \[\sum_{i=1}^m k_i = \sum_{i=1}^m \floor*{\frac{n}{p^i}}. \qh\]
\end{sol}

\begin{exercise} \label{ex0.2.9}
    Write a computer program to determine the greatest common divisor $(a, b)$ of two integers $a$ and $b$ and to express $(a, b)$ in the form $ax + by$ for some integers $x$ and $y$.
\end{exercise}

\begin{sol}
    Please see \verb|euclidean_algorithm.py| for the solution.
\end{sol}

\begin{exercise}\label{ex0.2.10}
    Prove for any given positive integer $N$ there exist only finitely many integers $n$ with $\phi(n) = N$ where $\phi$ denotes Euler's $\phi$-function. Conclude in particular that $\phi$ tends to infinity as $n$ tends to infinity.
\end{exercise}

\begin{sol}
    Fix $N \in \zp$, and consider $n \in \zp$ such that $\phi(n) = N$. Let $n = p_1^{\alpha_1} p_2^{\alpha_2} \cdots p_k^{\alpha_k}$ be the prime factorization of $n$. Then
    \[\phi(n) = \phi(p_1^{\alpha_1}) \phi(p_2^{\alpha_2}) \cdots \phi(p_k^{\alpha_k}) = \prod_{i=1}^k p_i^{\alpha_i - 1}(p_i - 1).\]
    Since $\phi(n) = N$, there are only finitely many possibilities for the primes $p_i$ and their exponents $\alpha_i$ that satisfy this equation. Hence, there are only finitely many integers $n$ such that $\phi(n) = N$.

    Assume, by way of contradiction, that $\phi$ does not tend to infinity as $n$ tends to infinity. Then there exists $M \in \zp$ such that $\phi(n) \leq M$ for infinitely many $n$. However, this contradicts the fact that for each fixed $N \in \zp$, there are only finitely many $n$ such that $\phi(n) = N$. Hence, $\phi$ must tend to infinity as $n$ tends to infinity.
\end{sol}

\begin{exercise}\label{ex0.2.11}
    Prove that if $d$ divides $n$ then $\phi(d)$ divides $\phi(n)$ where $\phi$ denotes Euler's $\phi$-function.
\end{exercise}

\begin{sol}
    Suppose $n$ has the prime factorization
    \[n = p_1^{\alpha_1} p_2^{\alpha_2} \cdots p_k^{\alpha_k},\]
    and let $d$ divide $n$. Then $d$ can be written as
    \[d = p_1^{\beta_1} p_2^{\beta_2} \cdots p_k^{\beta_k}\]
    for some integers $\beta_i$ with $0 \leq \beta_i \leq \alpha_i$ for all $1 \leq i \leq k$. Then
    \[\phi(n) = \prod_{i=1}^k p_i^{\alpha_i - 1}(p_i - 1) \quad \text{and} \quad \phi(d) = \prod_{i=1}^k p_i^{\beta_i - 1}(p_i - 1).\]
    It follows that $\phi(d) \mid \phi(n)$, since $\alpha_i - 1 \geq \beta_i - 1$ for all $1 \leq i \leq k$.
\end{sol}

\newpage

\subsection{\texorpdfstring{$\intmod$: The Integers Modulo $n$}{Z/nZ: The Integers Modulo n}}

\begin{exercise}\label{ex0.3.1}
    Write down explicitly all the elements in the residue classes of $\intmod[18]$.
\end{exercise}

\begin{sol}
    Note that $\bar x = \set{x + 18k \mid k \in \z}$. Explicitly, we have:
    \begin{align*}
        \bar 0 & = \set{0, 18, -18, 36, -36, \ldots} & \bar 1 & = \set{1, 19, -17, 37, -35, \ldots} & \bar 2 & = \set{2, 20, -16, 38, -34, \ldots} \\
        \bar 3 & = \set{3, 21, -15, 39, -33, \ldots} & \bar 4 & = \set{4, 22, -14, 40, -32, \ldots} & \bar 5 & = \set{5, 23, -13, 41, -31, \ldots} \\
        \bar 6 & = \set{6, 24, -12, 42, -30, \ldots} & \bar 7 & = \set{7, 25, -11, 43, -29, \ldots} & \bar 8 & = \set{8, 26, -10, 44, -28, \ldots} \\
        \bar 9 & = \set{9, 27, -9, 45, -27, \ldots} & \bar{10} & = \set{10, 28, -8, 46, -26, \ldots} & \bar{11} & = \set{11, 29, -7, 47, -25, \ldots} \\
        \bar{12} & = \set{12, 30, -6, 48, -24, \ldots} & \bar{13} & = \set{13, 31, -5, 49, -23, \ldots} & \bar{14} & = \set{14, 32, -4, 50, -22, \ldots} \\
        \bar{15} & = \set{15, 33, -3, 51, -21, \ldots} & \bar{16} & = \set{16, 34, -2, 52, -20, \ldots} & \bar{17} & = \set{17, 35, -1, 53, -19, \ldots} \qh
    \end{align*}
\end{sol}

\begin{exercise}\label{ex0.3.2}
    Prove that the distinct equivalence classes in $\intmod$ are precisely $\bar 0, \bar 1, \bar 2, \ldots, \bar{n - 1}$.

    \hint{Use the Division Algorithm.}
\end{exercise}

\begin{sol}
    Let $m \in \z$. By the Division Algorithm, there exists unique $q, r \in \z$ such that
    \[m = nq + r, \quad 0 \leq r < n.\]
    Then $m \equiv r \bmod n$, so $m \in \bar r$. By definition, $\bar r$ is one of $\bar 0, \bar 1, \ldots, \bar{n - 1}$. Thus, every equivalence class in $\intmod$ is one of $\bar 0, \bar 1, \ldots, \bar{n - 1}$. Moreover, they are all distinct. Otherwise, if $\bar x = \bar y$ for $0 \leq x, y < n$, then $x \equiv y \bmod n$, or $n \mid (x - y)$. However, since $-n < x - y < n$, we must have $x - y = 0$, or $x = y$.
\end{sol}

\begin{exercise}\label{ex0.3.3}
    Prove that if
    \[a = a_n 10^n + a_{n - 1}10^{n - 1} + \cdots + a_1 10 + a_0\]
    is any positive integer, then
    \[a \equiv (a_n + a_{n - 1} + \cdots + a_1 + a_0) \bmod 9.\]

    \hint{Use the fact that $10 \equiv 1 \bmod 9$.}

    \remark{Note that this is the usual arithmetic rule that the remainder after division by 9 is the same as the sum of the decimal digits mod 9—in particular an integer is divisible by 9 if and only if the sum of its digits is divisible by 9.}
\end{exercise}

\begin{sol} 
    Using the hint, then
    \begin{align*}
        a & \equiv (a_n10^n + a_{n - 1}10^{n - 1} + \cdots + a_110 + a_0) \bmod 9 \\
        & \equiv (a_n1^n + a_{n - 1}1^{n - 1} + \cdots + a_11 + a_0) \bmod 9 \\
        & = (a_n + a_{n - 1} + \cdots + a_1 + a_0) \bmod 9 \qh
    \end{align*}
\end{sol}

\begin{exercise}\label{ex0.3.4}
    Compute the remainder when $37^{100}$ is divided by $29$.
\end{exercise}

\begin{sol}
    One can calculate that $37^{64} \equiv 20 \bmod 29, 37^{32} \equiv 7 \bmod 29, 37^4 \equiv 7 \bmod 29$. Hence,
    \[37^{100} = 37^{64} \cdot 37^{32} \cdot 37^4 \equiv 20 \cdot 7 \cdot 7 \bmod 29 \equiv 23 \bmod 29. \qh\]
\end{sol}

\begin{exercise}\label{ex0.3.5}
    Compute the last two digits of $9^{1500}$.
\end{exercise}

\begin{sol}
    Recall that the last two digits of a number are given by its remainder when divided by 100. Note that $9^6 \equiv 41 \bmod 100$, while $9^4 \equiv 61 \bmod 100$. Hence, $9^{10} \equiv 1 \bmod 100$. Then
    \[9^{1500} = (9^{10})^{150} \equiv 1^{150} \bmod 100 \equiv 1 \bmod 100. \]
    Thus, the last two digits of $9^{1500}$ are $01$.
\end{sol}

\begin{exercise} \label{ex0.3.6}
    Prove that the square of the elements in $\intmod[4]$ are just $\bar 0$ and $\bar 1$.
\end{exercise}

\begin{sol}
    \begin{align*}
        0^2 & = 0 \equiv 0 \bmod 4 \\
        1^2 & = 1 \equiv 1 \bmod 4 \\
        2^2 & = 4 \equiv 0 \bmod 4 \\
        3^2 & = 9 \equiv 1 \bmod 4 \qh
    \end{align*}
\end{sol}

\begin{exercise} \label{ex0.3.7}
    Prove that for any integers $a$ and $b$ that $a^2 + b^2$ never leaves a remainder of $3$ when divided by $4$.

    \hint{Use \hyperref[ex0.3.6]{Exercise 0.3.6}.}
\end{exercise}

\begin{sol}
    By \hyperref[ex0.3.6]{Exercise 0.3.6}, $a^2$ and $b^2$ are both either $0 \bmod 4$ or $1 \bmod 4$, so their sum can be:
    \begin{align*}
        0 + 0 \equiv 0 \bmod 4 \\
        0 + 1 \equiv 1 \bmod 4 \\
        1 + 0 \equiv 1 \bmod 4 \\
        1 + 1 \equiv 2 \bmod 4
    \end{align*}
    In any of the sums, there is no remainder of $3$ when dividing by $4$.
\end{sol}

\begin{exercise}\label{ex0.3.8}
    Prove that the equation $a^2 + b^2 = 3c^2$ has no solutions in nonzero integers $a, b$ and $c$.
    
    \hint{Consider the equation mod 4 as in the previous two exercises and show that $a, b$ and $c$ would all have to be divisible by $2$. Then each of $a^2, b^2$ and $c^2$ has a factor of $4$ and by dividing through by $4$ show that there would be a smaller set of solutions to the original equation. Iterate to reach a contradiction.}
\end{exercise}

\begin{sol}
    We consider $a^2 + b^2 = 3c^2$ in mod 4. By \hyperref[ex0.3.7]{Exercise 0.3.7}, $a^2 + b^2$ is either 0, 1, or 2 mod 4. By \hyperref[ex0.3.6]{Exercise 0.3.6}, $c^2$ is either 0 or 1 mod 4, so that $3c^2$ is either 0 or 3 mod 4. Thus, both sides must be 0 mod 4.

    It is clear that $4 \mid 3 c^2$ implies $4 \mid c^2$, and hence $2 \mid c$. Observe that if $a$ and $b$ were both odd, then $a^2 + b^2 \equiv 2 \mod 4$. If one of $a$ or $b$ were odd and the other even, then $a^2 + b^2 \equiv 1 \mod 4$. Thus, both $a$ and $b$ must be even. Put $a = 2\alpha$, $b = 2\beta$, and $c = 2\gamma$ for nonzero $\alpha, \beta, \gamma \in \z$. Substituting into the original equation gives
    \[ (2\alpha)^2 + (2\beta)^2 = 3(2\gamma)^2.\]
    We may divide by 4 to obtain $\alpha^2 + \beta^2 = 3\gamma^2$, which is the same form as the original equation but with smaller positive integers. Repeating this process leads to an infinite descent, which is impossible. Therefore, there are no solutions in nonzero integers $a, b,$ and $c$.
\end{sol}

\begin{exercise}\label{ex0.3.9}
    Prove that the square of any odd integer always leaves a remainder of $1$ when divided by $8$.
\end{exercise}

\begin{sol}
    Put $a = 2k + 1$ for $k \in \z$. Then $a^2 = 4k(k + 1) + 1 \equiv 1 \bmod 8$ since $k(k + 1)$ is always even.
\end{sol}

\begin{specialexercise}\label{ex0.3.10}
    Prove that the number of elements of $\units$ is $\phi(n)$, where $\phi$ denotes the Euler $\phi$-function.
\end{specialexercise}

\begin{sol}
    It suffices to show that $\bar a \in \units$ if and only if $(a, n) = 1$.
    \begin{itemize}
        \item[\rightimp] If $\bar a \in \units$, then there exists $\bar b \in \intmod$ such that $\bar a \cdot \bar b = \bar 1$. Then $ab \equiv 1 \bmod n$ implies the existence of $k \in \z$ such that $ab - 1 = kn$, or $ab - kn = 1$. Thus, $(a, n) = 1$.
        \item[\leftimp] By $(a, n) = 1$, there exists $x, y \in \z$ such that $ax + ny = 1$, or $ax = 1 - ny$. We then have $ax \equiv 1 \bmod n$, so that $\bar a \cdot \bar x = \bar 1$. Thus, $\bar a \in \units$. \qh
    \end{itemize}
\end{sol}

\begin{exercise}\label{ex0.3.11}
    Prove that if $\bar a, \bar b \in \units$, then $\bar a \cdot \bar b \in \units$.
\end{exercise}

\begin{sol}
    Since $\bar a, \bar b \in \units$, there exist $\bar c, \bar d \in \intmod$ such that $\bar a \cdot \bar c = \bar b \cdot \bar d = \bar 1$. Then
    \[(\bar a \cdot \bar b)(\bar c \cdot \bar d) = (\bar a \cdot \bar c)(\bar b \cdot \bar d) = \bar 1 \cdot \bar 1 = \bar 1,\]
    so that $\bar a \cdot \bar b \in \units$.
\end{sol}

\begin{exercise}\label{ex0.3.12}
    Let $n \in \z, n > 1$ and let $a \in \z$ with $1 \leq a \leq n$. Prove if $a$ and $n$ are not relatively prime, there exists an integer $b$ with $1 \leq b < n$ such that $ab \equiv 0 \bmod n$ and deduce that there cannot be an integer $c$ such that $ac \equiv 1 \bmod n$.
\end{exercise}

\begin{sol}
    Let $(a, n) = d > 1$. Then $a/d, n/d \in \z$ and $1 \leq a/d < n/d$. Put $b = n/d$, so that $1 \leq b < n$ and
    \[ab = a(n/d) = (a/d)n \equiv 0 \bmod n.\]
    If there existed some $c \in \z$ such that $ac \equiv 1 \bmod n$, then $(ac)b \equiv b \bmod n$. However, $ab \equiv 0 \bmod n$, so that $0 \equiv b \bmod n$, which is impossible since $1 \leq b < n$.
\end{sol}

\begin{exercise}\label{ex0.3.13}
    Let $n \in \z, n > 1$ and let $a \in \z$ with $1 \leq a \leq n$. Prove that if $a$ and $n$ are relatively prime then there is an integer $c$ such that $ac \equiv 1 \bmod n$.

    \hint{Use the fact that the GCD of two integers is a $\z$-linear combination of the integers.}
\end{exercise}

\begin{sol}
    Since $(a, n) = 1$, there exists $c, x \in \z$ such that $ac + nx = 1$. Then $ac \equiv 1 \bmod n$.
\end{sol}

\begin{exercise}\label{ex0.3.14}
    Conclude from the previous two exercises that $\units$ is the set of elements $\bar a$ of $\intmod$ with $(a, n) = 1$ and hence prove Proposition 4. Verify this directly in the case $n = 12$.
\end{exercise}

\begin{sol}
    The previous two exercises show that $a$ and $n$ are relatively prime if and only if there exists $b$ such that $ab \equiv 1 \bmod n$, which is exactly the proposition. For $n = 12$, the elements $1, 5, 7, 11$ are relatively prime to $12$, so that $\units[11] = \{\bar 1, \bar 5, \bar 7, \bar{11}\} $ whose inverses are $\bar 1, \bar 7, \bar 5, \bar{11}$ respectively. 
\end{sol}

\begin{exercise}\label{ex0.3.15}
    For each of the following pairs of integers $a$ and $n$, show that $a$ is relatively prime to $n$ and determine the multiplicative inverse of $\bar a$ in $\intmod$.
    \begin{subproblems}
        \item $a = 13, n = 20$
        \item $a = 69, n = 89$
        \item $a = 1891, n = 3797$
        \item $a = 6003722857, n = 77695236973$
    \end{subproblems}
\end{exercise}

\begin{sol}
    Refer to \hyperref[ex0.2.1]{Exercise 0.2.1} to see how the Euclidean Algorithm is used to find $x, y \in \z$ such that $ax + ny = 1$. From this, it follows that $ax \equiv 1 \bmod n$, so that $\bar x$ is the multiplicative inverse of $\bar a$ in $\intmod$.
    \begin{subproblems}
        \item $-3(13) + 2(20) = 1 \implies -3 \cdot 13 \equiv 1 \bmod 20$, so the inverse is $\bar{-3} = \bar{17}$.
        \item $40(69) - 31(89) = 1 \implies 40 \cdot 69 \equiv 1 \bmod 89$, so the inverse is $\bar{40}$.
        \item $253(1891) - 126(3797) = 1 \implies 253 \cdot 1891 \equiv 1 \bmod 3797$, so the inverse is $\bar{253}$.
        \item $17n - 220a = 1 \implies -220a \equiv 1 \bmod n \implies 77695237193a \equiv 1 \bmod n$, so the inverse is $\bar{77695237193}$. \qh
    \end{subproblems}
\end{sol}

\begin{exercise}\label{ex0.3.16}
    Write a computer program to add and multiply mod $n$, for any $n$ given as input. The output of these operations should be the least residues of the sums and products of two integers. Also include the feature that if $(a, n) = 1$, an integer $c$ between $1$ and $n - 1$ such that $\bar a \cdot \bar c = \bar 1$ may be printed on request. 
    
    \remark{Your program should not, of course, simply quote ``mod'' functions already built into many systems.}
\end{exercise}

\begin{sol}
    Please see \verb|modular_arithmetic.py| for the solution.
\end{sol}